\section{Stability}\label{sec:stability}

The Stability axis asks when declared future behavior is already
controlled by present data, and when a residual must be recorded
instead of ignored.  In the language of \Cref{sec:apparatus}, the
question is whether the current quotient carries the probes the
package has declared, or whether closure-stability fails under a
structural perturbation.  A perturbation may be a new future probe,
a layer-dissolving move at a membrane, a strict extension of a
quotient, or an unpriced residual.  The axis records the normal
forms that decide which of these outcomes is legitimate.
In the running example of \Cref{sec:intro}, the value of \(F\) is a
declared future probe that the current quotient \(q\) does not carry:
\(q\) identifies \(1\) and \(2\) while \(F\) separates them, so the
split pair must be resolved or recorded as a residual.

The chapter has four laws.  No-Needles Stability is the first theorem of the chapter: when the
currency, residual and endpoint budgets hold together with the stated block
identities, every transported dissolving direction is bounded by its budget,
so no needle exists.  Sufficiency Closure then gives the
factorization criterion for declared futures: the current quotient is
sufficient exactly when every declared future probe descends through
it.  Strict Extension records when genuine layer growth has occurred,
namely when the old quotient is retained and a new quotient splits at
least one old fiber.  No Unstatused Residual closes the axis: under the
residual-coverage rule and agreement of statuses at a common locus, every
native residual carries a well-defined typed status, and neither hypothesis
can be dropped.

No-Needles Stability is the layer-agnostic form of the layer-dissolving
membrane theorem of \citet{Tsiokos2026NeedleKiller}, which
\citet{Tsiokos2026NavierStokes} applies to the Navier--Stokes
equations.  Here a membrane means the
layer-dissolving boundary of a closure: the place where a distinction
between ``inside'' and ``outside'' the layer can cease to be
meaningful.  Strict Extension also has descriptive substrate context,
but it is a direct law here: the Schur-calculus intuition comes from
\PaperAR\ \citep{Tsiokos2026XiCompanion}, where residuals compose
along refinement chains.  The Stability axis can be read
independently of the Access axis; its only prerequisite is the
notation of \Cref{sec:apparatus}.

\subsection{No-Needles Stability [\texorpdfstring{\Fref{F6}}{F6}]}\label{sec:stability:no-needles}

The setup is a formed stability package: a closure carrier equipped
with the inherited closure apparatus, quotient discipline, audit
ledger, and a layer-dissolving membrane.  The membrane is the
boundary across which a native layer distinction is allowed to
dissolve only when the package has recorded the appropriate
apparatus footprint.  A \emph{needle} is a transported dissolving direction whose effect exceeds
its declared budget. In the matrix data of the theorem below, it is a vector
\(v\in\mathbb R^d\) with
\(v^{\!\top}S\,K_{DD}\,S^{\!\top}v>v^{\!\top}\Theta_D v\) (the block \(K_{DD}\), the selector \(S\) and the bound \(\Theta_D\in\mathbb R^{d\times d}\) are introduced in the paragraphs \emph{Notation} and \emph{Bounds and order} below).

\paragraph{Notation.} The theorem itself concerns four blocks
\(K_{LL},K_{LD},K_{DL},K_{DD}\) and a matrix \(K_{LL}^{\dagger}\) subject to the
identities displayed in it; the next paragraph describes their
intended origin. The data of the theorem are the matrices displayed in it.
\paragraph{Intended origin (not used by the theorem).} In the intended origin, which the theorem does not use, the currency
block \(C\) is an \(e\times e\) matrix on the exterior carrier, and \(L\)
and \(D\) are \(y\times e\) and \(z\times e\) blocks for the layer-bridge
and dissolving directions. Let \(C^{+}\) be any symmetric \(e\times e\) matrix (for instance the
Moore--Penrose pseudoinverse of a symmetric \(C\); the theorem uses \(C\) only through \(C^{+}\), and uses the blocks only
through the stated identities); the currency pieces are \(K_{LL}=L\,C^{+}L^{\!\top}\),
\(K_{LD}=L\,C^{+}D^{\!\top}\), \(K_{DL}=D\,C^{+}L^{\!\top}\) and
\(K_{DD}=D\,C^{+}D^{\!\top}\), and symmetry of \(C^{+}\) gives
\(K_{DL}=K_{LD}^{\!\top}\). Let \(K_{LL}^{\dagger}\) be any symmetric \(y\times y\)
matrix with \(K_{LL}K_{LL}^{\dagger}K_{LD}=K_{LD}\) (for instance the
Moore--Penrose pseudoinverse of \(K_{LL}\) when \(C^{+}\succeq0\)). The declared native explanation
\(\widehat E=K_{DL}K_{LL}^{\dagger}\) predicts the dissolving block from the
layer-bridge block, and the adequacy residual
\(\Xi=K_{DD}-K_{DL}K_{LL}^{\dagger}K_{LD}\), a generalized Schur complement,
is the algebraic remainder of that explanation. When \(C^{+}\succeq0\), the
block matrix \(\bigl[\begin{smallmatrix}K_{LL}&K_{LD}\\K_{DL}&K_{DD}\end{smallmatrix}\bigr]\)
is positive semidefinite; then \(\Xi\succeq0\), and for every \(z\times y\)
matrix \(E\),
\(K_{DD}-EK_{LD}-K_{DL}E^{\!\top}+EK_{LL}E^{\!\top}=\Xi+(E-\widehat E)K_{LL}(E-\widehat E)^{\!\top}\succeq\Xi\),
so \(\widehat E\) is the Loewner-optimal explanation. For arbitrary
symmetric \(C^{+}\), neither positivity nor optimality is asserted; the endpoint proof uses the block identities and
budgets alone.
\paragraph{Bounds and order.} \(\Theta_Y\in\mathbb R^{y\times y}\), \(\Omega_Z\in\mathbb R^{z\times z}\) and
\(\Theta_D\in\mathbb R^{d\times d}\) are declared bounds; no construction from the blocks \(L\) or \(D\) is required,
\(\preceq\) is the Loewner preorder of quadratic forms
(\(A\preceq B\) when \(v^{\!\top}(B-A)v\ge0\) for every vector \(v\); on
symmetric matrices this is the Loewner order, and in general it compares
symmetric parts), and \(S\) is any \(d\times z\) matrix, for instance
one selecting \(d\) of the dissolving directions.
In \citet{Tsiokos2026NeedleKiller} the audit form \(C\) is a positive
semidefinite operator on a finite-dimensional real or complex space, and
\(C^{+}\) and \(K_{LL}^{\dagger}\) are Moore--Penrose pseudoinverses.  The
proof below uses only the displayed matrix identities and budgets; the
package and its membrane name the setting in which they are declared.

\begin{theorem}[No-Needles Stability]\label{thm:F6}
Suppose given the following data (in the intended reading, declared by a
formed stability package \(\mathcal P\), whose other properties, including its
membrane, are not used):
\begin{itemize}
  \item block matrices \(K_{LL}\in\mathbb R^{y\times y}\),
  \(K_{LD}\in\mathbb R^{y\times z}\), \(K_{DL}\in\mathbb R^{z\times y}\) and
  \(K_{DD}\in\mathbb R^{z\times z}\) (in the intended reading, the currency pieces
  formed from a symmetric \(C^{+}\) and the blocks \(L\), \(D\) as in the paragraph \emph{Intended origin}) and a matrix \(K_{LL}^{\dagger}\in\mathbb R^{y\times y}\);
  \item the symmetry conditions \(K_{DL}=K_{LD}^{\!\top}\) (automatic when the
  pieces are formed from a symmetric \(C^{+}\)) and
  \(K_{LL}^{\dagger}=(K_{LL}^{\dagger})^{\!\top}\), together with the
  projection identity
  \[
    K_{LL}\,K_{LL}^{\dagger}\,K_{LD}=K_{LD};
  \]
  \item currency and residual budgets
  \[
    K_{LL}\preceq \Theta_Y,
    \qquad
    \Xi\preceq \Omega_Z,
  \]
  where \(\Xi:=K_{DD}-K_{DL}K_{LL}^{\dagger}K_{LD}\) is the adequacy residual;
  \item a transport selector \(S\) of size \(d\times z\) satisfying
  the endpoint budget
  \[
    S\bigl(\widehat E\,\Theta_Y\,\widehat E^{\!\top}
      +\Omega_Z\bigr)S^{\!\top}\preceq\Theta_D,
  \]
  where \(\widehat E:=K_{DL}K_{LL}^{\dagger}\) is the declared native
  explanation.
\end{itemize}
Then the transported dissolving block satisfies the Loewner endpoint
\[
  S\,K_{DD}\,S^{\!\top}\preceq\Theta_D .
\]
Equivalently, no needle exists: every transported dissolving direction is
bounded by its declared budget.
\end{theorem}
\begin{proof}
Assume the symmetry, projection, currency-budget, residual-budget,
and transport-budget hypotheses displayed in the theorem.  The
projection identity
\[
  K_{LL}\,K_{LL}^{\dagger}\,K_{LD}=K_{LD}
\]
gives the layer-dissolving decomposition of the substrate master theorem.
Indeed, with \(\widehat E:=K_{DL}\,K_{LL}^{\dagger}\), the symmetry conditions
give \(\widehat E^{\!\top}=K_{LL}^{\dagger}K_{LD}\), so
\[
  \widehat E\,K_{LL}\,\widehat E^{\!\top}
  =K_{DL}K_{LL}^{\dagger}\bigl(K_{LL}K_{LL}^{\dagger}K_{LD}\bigr)
  =K_{DL}K_{LL}^{\dagger}K_{LD},
\]
and by the definition of \(\Xi\),
\[
  K_{DD}
  =
  \widehat E\,K_{LL}\,\widehat E^{\!\top}
  +\Xi.
\]
Here \(\widehat E\) is the declared native explanation and \(\Xi\) is
the adequacy residual.  Applying Loewner monotonicity under the
declared bounds \(K_{LL}\preceq\Theta_Y\) and
\(\Xi\preceq\Omega_Z\) yields
\[
  K_{DD}
  \preceq
  \widehat E\,\Theta_Y\,\widehat E^{\!\top}+\Omega_Z .
\]
Here the first step also uses that congruence preserves the order: for any
matrix \(A\) and any \(v\),
\(v^{\!\top}A(Y-X)A^{\!\top}v=(A^{\!\top}v)^{\!\top}(Y-X)(A^{\!\top}v)\), so
\(X\preceq Y\) implies \(AXA^{\!\top}\preceq AYA^{\!\top}\), applied with
\(A=\widehat E\). Conjugation by \(S\) preserves the Loewner order in the same
way, so
\[
  S\,K_{DD}\,S^{\!\top}
  \preceq
  S\bigl(\widehat E\,\Theta_Y\,\widehat E^{\!\top}
    +\Omega_Z\bigr)S^{\!\top}.
\]
The declared transport-side endpoint budget gives
\[
  S\bigl(\widehat E\,\Theta_Y\,\widehat E^{\!\top}
    +\Omega_Z\bigr)S^{\!\top}
  \preceq \Theta_D .
\]
Transitivity of \(\preceq\) therefore gives the endpoint
\[
  S\,K_{DD}\,S^{\!\top}\preceq\Theta_D .
\]

\end{proof}

The proof uses only the arithmetic of an ordered field and the order
on quadratic forms, so the theorem holds verbatim with \(\mathbb R\)
replaced by any ordered field, for instance \(\mathbb Q\).  The complex
instances of \citet{Tsiokos2026NeedleKiller} need the analogous
statement with conjugate transposes and Hermitian forms, which is not
stated here.

\paragraph{What the sources supply.} The layer-dissolving membrane
theorem of \citet{Tsiokos2026NeedleKiller} applies the same argument at
every accepted stage of a refinement, with positive semidefinite audit
forms and Moore--Penrose pseudoinverses, and identifies the endpoint
with a bound on delivered capacity; \citet{Tsiokos2026NavierStokes}
applies it to shell readouts of the Navier--Stokes equations.
\Cref{thm:F6} isolates the matrix inequality common to these
instances.

\paragraph{Nonclaims.} The theorem does not re-prove the Needle Killer or Navier--Stokes
substrate instances, does not classify every possible
layer-dissolving move, and does not assert that the substrate carrier
is unique.  It records the common stability normal form: under the
declared budgets, the transported dissolving block is dominated in the Loewner
order by its endpoint budget.

\subsection*{Layer instantiations}\label{layer-inst:F6}

\begingroup\small
\begin{description}
\item[\emph{Schur-complement positivity criterion (linear algebra).}] \textit{Analogy.}
A symmetric block matrix
\(\bigl[\begin{smallmatrix}A&B\\B^{\!\top}&C\end{smallmatrix}\bigr]\)
is positive semidefinite iff \(A\succeq0\), \((I-AA^{+})B=0\), and
\(C-B^{\!\top}A^{+}B\succeq0\).  The criterion is not a case of \Cref{thm:F6}; it is the positivity fact stated in the paragraph \emph{Intended origin}, with \(A,B,C,A^{+}\) in the roles of \(K_{LL},K_{LD},K_{DD},K_{LL}^{\dagger}\): \((I-AA^{+})B=0\) is the projection identity \(K_{LL}K_{LL}^{\dagger}K_{LD}=K_{LD}\), and the Schur complement \(C-B^{\!\top}A^{+}B\) is the adequacy residual \(\Xi\), whose bound \(\Xi\preceq\Omega_Z\) is a hypothesis of the theorem; the transported endpoint is \(SK_{DD}S^{\!\top}\) \citep{Zhang2005SchurApplications}.

\end{description}
\endgroup


\subsection{Sufficiency Closure [\texorpdfstring{\Fref{F7}}{F7}]}\label{sec:stability:sufficiency-closure}

Fix a closure package with history carrier \(H\), current quotient
map \(q:H\to Q\), and a declared family \(\mathcal F\) of future
probes.  The family is the apparatus's announcement of which future
observables will be tested at this interface; it is not a claim of
sufficiency for every possible future question.

Write \(\Phi:H\to\prod_{f\in\mathcal F}F_f\), \(\Phi(h)=(f(h))_{f\in\mathcal F}\),
for the bundled future signature, and call \(q\) \emph{sufficient} for
\(\mathcal F\) when \(q(h)=q(h')\) implies \(f(h)=f(h')\) for every
\(f\in\mathcal F\).

\begin{theorem}[Sufficiency Closure]\label{thm:F7}
Assume \(q\) is surjective. The current quotient \(q\) is sufficient
for the declared future family \(\mathcal F\) if and only if every declared
future probe \(f\in\mathcal F\) descends through \(q\), that is, for every
\(f:H\to F_f\) in \(\mathcal F\) there is a map
\(\overline f:Q\to F_f\) such that
\[
  f=\overline f\circ q ,
\]
equivalently \(\Phi=\overline\Phi\circ q\) for some
\(\overline\Phi:Q\to\prod_{f\in\mathcal F}F_f\), and if and only if no pair \((h,h')\) has \(q(h)=q(h')\) and
\(\Phi(h)\ne\Phi(h')\). Moreover exactly one of the following four statuses
holds: \emph{exact closure} (\(q\) is sufficient and \(\Phi\) determines \(q\): \(\Phi(h)=\Phi(h')\Rightarrow q(h)=q(h')\) for all \(h,h'\in H\)); \emph{redundant closure} (\(q\) sufficient, \(\Phi\) does not determine \(q\)); \emph{predictive non-closure} (\(\Phi\) determines \(q\), \(q\) not sufficient); and \emph{incomparable} (neither).
\end{theorem}
\begin{proof}
Assume first that every declared future probe descends through
\(q\).  Then each future value \(f(h)\) is determined by the current
class \(q(h)\), so two histories with the same current class have the
same value under every declared probe; this is sufficiency.

Conversely, suppose \(q\) is sufficient for the declared future family.
For a probe \(f:H\to F_f\) and a class \(x\in Q\), choose \(h\) with
\(q(h)=x\), which exists because \(q\) is surjective, and set
\(\overline f(x)=f(h)\). Sufficiency makes this independent of the chosen
\(h\): if \(q(h)=q(h')\), then \(f(h)=f(h')\). Therefore \(\overline f\) is
well-defined and satisfies \(f=\overline f\circ q\). The bundled form follows
by taking \(\overline\Phi=(\overline f)_{f\in\mathcal F}\); conversely
\(\overline f=\mathrm{pr}_f\circ\overline\Phi\).

Sufficiency says that no pair in a common current fiber is separated by a
probe, which is the statement that no pair has \(q(h)=q(h')\) and
\(\Phi(h)\ne\Phi(h')\). Finally, the four statuses are the four
combinations of two Boolean conditions, sufficiency and the implication
\(\Phi(h)=\Phi(h')\Rightarrow q(h)=q(h')\); every configuration satisfies
exactly one combination.
\end{proof}

\paragraph{Nonclaims.} The theorem does not enumerate the
declared future probes; that choice belongs to the apparatus.  It
also does not assert sufficiency for undeclared probes, which lie
outside the current audit scope.

\subsection*{Layer instantiations}\label{layer-inst:F7}

\begingroup\small
\begin{description}
\item[\emph{Information bottleneck (information theory / machine
learning).}] \textit{Special case.}
For input \(X\) and a declared prediction target \(Y\), the
Tishby--Pereira--Bialek information bottleneck trades the compression cost \(I(X;T)\) against the relevance
\(I(T;Y)\) over stochastic encoders; its limiting case, with a
deterministic encoder \(T=q(X)\), \(q\) corestricted to its image as \Fref{F7} requires, minimises \(I(X;T)\) under the
sufficiency constraint \(I(T;Y)=I(X;Y)\).  For a finite input alphabet
with \(P_X(x)>0\) for every \(x\), taken as \(H\), the constraint holds exactly
when \(P(Y\mid X=x)=P(Y\mid T=q(x))\) for every \(x\), which is the \Fref{F7}
factorization for \(\mathcal F=\{x\mapsto P(Y\mid X=x)\}\).  Over-compressing the
encoder breaks the factorization: the \Fref{F7} obstruction is nonempty
exactly when the gap \(I(X;Y)-I(T;Y)\) is positive
\citep{TishbyPereiraBialek1999InformationBottleneck}.

\item[\emph{Predictive state representations (reinforcement
learning).}] \textit{Special case.}
A PSR summarizes the history of a partially observable dynamical
system by the vector \(p(h)\), corestricted to its image as \Fref{F7} requires, of probabilities of a declared family of
core future tests.  The Singh--James--Rudary sufficiency theorem
states that every declared future test factors linearly through
\(p(h)\): \(\Pr(t\mid h)=m_t\cdot p(h)\).  The summary is sufficient
for the declared test family while remaining hidden-state-free, and a
summary built from fewer tests than the system's linear dimension does not make
every test a linear function of it
\citep{SinghJamesRudary2004PredictiveStateRepresentations}.

\item[\emph{Query-complete materialized views (databases).}] \textit{Special case.}
For a base schema \(R\) and a declared workload \(\mathcal Q\) of
queries, a set of materialized views \(V\) \emph{determines} \(Q\in\mathcal Q\)
when \(V(I)=V(I')\) implies \(Q(I)=Q(I')\) for all instances \(I,I'\). By
\Fref{F7}, with the view map corestricted to its image, this holds iff
\(Q=Q'\circ V\) for some map \(Q'\) on view extents; the view materialization
plays the summary and \(Q'\) the lower map. A rewriting of \(Q'\) in a query
language is a stronger requirement: conjunctive views can determine a
conjunctive query that has no conjunctive rewriting \citep{NashSegoufinVianu2010Views}.  The
sufficiency is workload-relative: a query needing an attribute
projected away by \(V\) does not factor through the views
\citep{Halevy2001AnsweringQueriesUsingViewsSurvey}.

\item[\emph{Quantum sufficient subalgebras (quantum information).}] \textit{Analogy.}
For a parametric family of density operators
\(\{\rho_\theta:\theta\in\Theta\}\) and a CPTP quantum channel \(T\),
Petz's theorem says \(T\) is sufficient for the family iff there is a
recovery channel \(R\) (the Petz recovery map) with
\(\rho_\theta=R(T(\rho_\theta))\) for every \(\theta\).  Every declared
measurement statistic \(\operatorname{Tr}(\rho_\theta M_i)\) then
factors through \(T\), and the obstruction is strict contraction of
relative entropy by \(T\), which prevents recovery
\citep{Petz1986SufficientSubalgebras}.

\item[\emph{Noninterference for low observers.}] \textit{Special case.}
For a deterministic program \(P:H\to H\) on states with high-security (secret) and low-security (public) components, take the current quotient to be the low projection \(q:H\to H{\restriction}L\) of the initial state and the declared probes to be \(h\mapsto o(P(h))\) for the declared low observers \(o\) of the final state; \(P\) satisfies noninterference iff every such probe factors through \(q\).  Side-channel attacks --- timing, cache,
fault --- appear as family-extension bridges: the declared
low-observer family is enlarged to include an observer for which the
low-projection is no longer sufficient
\citep{GoguenMeseguer1982SecurityPolicies,SabelfeldMyers2003LanguageBasedInformationFlow}.
\end{description}
\endgroup

\subsection{Strict Extension [\texorpdfstring{\Fref{F8}}{F8}]}\label{sec:stability:strict-extension}

Let \(H\) be the carrier under discussion, let
\[
  q_{\mathrm{old}}:H\to Q_{\mathrm{old}}
\]
be the old quotient, and let
\[
  q_{\mathrm{new}}:H\to Q_{\mathrm{new}}
\]
be a candidate new quotient.  A new presentation is not
automatically a new layer: it may be a relabeling, a coarsening, or
an incomparable rewrite.  The Schur-calculus intuition from
\PaperAR\ \citep{Tsiokos2026XiCompanion} is that residuals compose
along refinement chains, so genuine layer growth preserves the old
quotient while accounting for the new split.

An extension is \emph{strict} when it retains the old quotient, so that the
old classes can be read off the new ones, and adds genuine growth, so that
it is not a relabeling of the old quotient. Precisely, the extension from
\(q_{\mathrm{old}}\) to \(q_{\mathrm{new}}\) \emph{retains} the old quotient
when \(q_{\mathrm{old}}=\rho\circ q_{\mathrm{new}}\) for some map \(\rho\), the
new quotient is \emph{definable from the old} when
\(q_{\mathrm{new}}=\lambda\circ q_{\mathrm{old}}\) for some map \(\lambda\), and
the extension is strict when it retains the old quotient and the new one is
not definable from the old. The lemma turns the second condition into a
split fiber and places strictness among the four possible comparisons.

\begin{lemma}[Strict Extension]\label{lem:F8}
Let \(q_{\mathrm{old}}\) be surjective. The extension from \(q_{\mathrm{old}}\) to \(q_{\mathrm{new}}\) is
strict if and only if there is a retention map
\[
  \rho:Q_{\mathrm{new}}\to Q_{\mathrm{old}},
  \qquad
  q_{\mathrm{old}}=\rho\circ q_{\mathrm{new}},
\]
and there are histories \(h,h'\in H\) such that
\[
  q_{\mathrm{old}}(h)=q_{\mathrm{old}}(h')
  \quad\text{but}\quad
  q_{\mathrm{new}}(h)\ne q_{\mathrm{new}}(h').
\]
Moreover, for any \(q_{\mathrm{old}}\) and \(q_{\mathrm{new}}\), surjective or not, exactly one of
four comparisons holds: \emph{equivalent} (retains and definable),
\emph{strict} (retains and not definable), \emph{coarsening} (definable and
not retains) or \emph{incomparable} (neither).
\end{lemma}
\begin{proof}
Assume the extension is strict.  By the retention part of strictness,
there is a map \(\rho:Q_{\mathrm{new}}\to Q_{\mathrm{old}}\) with
\(q_{\mathrm{old}}=\rho\circ q_{\mathrm{new}}\).  If no old fiber were
split, then \(q_{\mathrm{new}}\) would be constant on each fiber of
\(q_{\mathrm{old}}\); since \(q_{\mathrm{old}}\) is surjective, every
\(y\in Q_{\mathrm{old}}\) is \(q_{\mathrm{old}}(h)\) for some \(h\), and
\(\lambda(y):=q_{\mathrm{new}}(h)\) would be a well-defined map with
\(q_{\mathrm{new}}=\lambda\circ q_{\mathrm{old}}\), contradicting
non-definability.  So there are \(h,h'\in H\) with
\(q_{\mathrm{old}}(h)=q_{\mathrm{old}}(h')\) but
\(q_{\mathrm{new}}(h)\ne q_{\mathrm{new}}(h')\).

Conversely, suppose such a retention map and split old fiber are
given.  The equation \(q_{\mathrm{old}}=\rho\circ q_{\mathrm{new}}\)
says that every new class lies inside an old class, so the candidate
retains the old quotient rather than replacing it.  The displayed
pair \(h,h'\) shows that at least one old class has been split by
the new quotient, so no \(\lambda\) satisfies
\(q_{\mathrm{new}}=\lambda\circ q_{\mathrm{old}}\): it would give
\(q_{\mathrm{new}}(h)=\lambda(q_{\mathrm{old}}(h))=\lambda(q_{\mathrm{old}}(h'))=q_{\mathrm{new}}(h')\).
Thus the extension is strict.  The four comparisons are the four
combinations of the two conditions, so exactly one of them holds.
\end{proof}

\paragraph{Nonclaims.} In the strict case the lemma supplies a
retention map \(\rho\) and a split pair \((h,h')\), which a finite audit
can check.  It does not re-derive the Schur calculus,
and it does not classify every kind of layer growth.  It only gives
the quotient test for genuine retentive growth.

\subsection*{Layer instantiations}\label{layer-inst:F8}

\begingroup\small
\begin{description}
\item[\emph{Ring species (evolutionary biology).}] \textit{Analogy.}
For a ring-distributed population complex (the greenish-warbler ring around the Tibetan Plateau, the
\emph{Ensatina} salamander complex around California's Central
Valley), the local interbreeding equivalence
\(q_{\mathrm{old}}\) identifies adjacent populations that interbreed
in the wild --- yielding a single ring-wide class.  Reproductive compatibility is not
transitive around the ring, so it is not itself an equivalence; the
reproductive partition \(q_{\mathrm{new}}\) refines
\(q_{\mathrm{old}}\) by assigning the two reproductively isolated
terminal forms to distinct classes at a declared boundary along the
chain.  The
retention map is the inclusion of new classes into the single old
class, and the witness split is the closure pair, reproductively
isolated yet connected through the chain of locally interbreeding
intermediates \citep{IrwinBenschPrice2001SpeciationInARing}.

\item[\emph{Refinement mappings in formal verification (computer
science).}] \textit{Analogy.}
For a high-level specification \(S_2\) and a candidate implementation
\(S_1\), the Abadi--Lamport refinement-mapping theorem states that, when
\(S_1\) is machine-closed and \(S_2\) has finite invisible
nondeterminism and is internally continuous, \(S_1\) implements
\(S_2\) if and only if, possibly after adjoining history and prophecy
variables, there exists a state-level map
\(f:\mathrm{States}(S_1)\to\mathrm{States}(S_2)\) sending each \(S_1\)
step to a legal \(S_2\) step or a stuttering step and preserving
liveness.  The map \(f\) is the \Fref{F8} retention map;
the witness split is two implementation behaviors with identical
\(S_2\)-observable trace but distinct internal scheduling or
auxiliary-variable values
\citep{AbadiLamport1991RefinementMappings}.

\end{description}
\endgroup


\subsection{No Unstatused Residual [\texorpdfstring{\Fref{F9}}{F9}]}\label{sec:stability:no-unstatused-residual}

A formed stability package has a native residual registry: the
defects, gaps, budget failures, split pairs, critical pairs, or
endpoint failures that its own claims make relevant.  Each native
residual is not merely a number or a warning flag.  It comes with a
record saying which judgment family is responsible for it and which
status classifier applies.

\paragraph{Reading the statement.}
\(\mathcal S_{\mathrm{typed}}\) is the set of typed statuses the package
may assign: a status value, such as visible, hidden with a record, outside scope,
blocked, priced or accepted, tagged by the judgment family that assigns it.
The package keeps a set \(\mathrm{Rec}\) of residual records; each record
\(\rho\) names a residual \(\operatorname{loc}(\rho)\in R\), its locus, and
carries a typed status \(\operatorname{st}(\rho)\in\mathcal S_{\mathrm{typed}}\).
A residual \(r\) has a typed status when it is the locus of some record, and
it is \emph{unstatused}, \(\operatorname{UnstatusedResidual}_{\mathcal P}(r)\),
when it is the locus of no record.

\begin{theorem}[No Unstatused Residual]\label{thm:F9}
Let \(R\) be a residual type and let \(\mathcal P\) be a stability package over
\(R\) (in the intended reading, a formed one) with residual records \(\mathrm{Rec}\), locus map
\(\operatorname{loc}:\mathrm{Rec}\to R\) and status map
\(\operatorname{st}:\mathrm{Rec}\to\mathcal S_{\mathrm{typed}}\). Assume
the residual-coverage rule (RC) of \Cref{sec:apparatus}, so every residual is
the locus of a record, and assume records with the same locus carry the same
status: \(\operatorname{loc}(\rho)=\operatorname{loc}(\rho')\) implies
\(\operatorname{st}(\rho)=\operatorname{st}(\rho')\). (This holds in particular
when records are unique at their locus.) Then, for every \(r\in R\):
\begin{enumerate}[label=(\roman*)]
\item \(r\) has a typed status;
\item all records with locus \(r\) carry the same typed status, so
  \(\operatorname{st}_{\mathcal P}(r):=\operatorname{st}(\rho)\) for any
  record \(\rho\) with locus \(r\) defines a map
  \(\operatorname{st}_{\mathcal P}:R\to\mathcal S_{\mathrm{typed}}\);
\item \(\neg\operatorname{UnstatusedResidual}_{\mathcal P}(r)\).
\end{enumerate}
Clauses (i) and (iii) restate (RC) at \(r\); clause (ii) uses the agreement
of statuses at a common locus. Both hypotheses are needed: under (RC), a map
\(\operatorname{st}_{\mathcal P}:R\to\mathcal S_{\mathrm{typed}}\) with
\(\operatorname{st}_{\mathcal P}(\operatorname{loc}\rho)=\operatorname{st}(\rho)\)
for every record exists if and only if records with a common locus carry the
same status; and (RC) is equivalent to the absence of unstatused residuals.
\end{theorem}
\begin{proof}
Let \(r\in R\). By (RC) there is a record \(\rho\) with
\(\operatorname{loc}(\rho)=r\), and \(\operatorname{st}(\rho)\) is a typed
status of \(r\); this is (i). If \(\rho'\) is another record with
\(\operatorname{loc}(\rho')=r\), the agreement hypothesis gives
\(\operatorname{st}(\rho')=\operatorname{st}(\rho)\); this
is (ii), and \(\operatorname{st}_{\mathcal P}(r)\) does not depend on the
record chosen. Finally, \(\operatorname{UnstatusedResidual}_{\mathcal P}(r)\)
says that \(r\) is the locus of no record, which contradicts (RC); this is
(iii). For the converses, if a map \(\operatorname{st}_{\mathcal P}\) with
\(\operatorname{st}_{\mathcal P}(\operatorname{loc}\rho)=\operatorname{st}(\rho)\)
exists and \(\operatorname{loc}(\rho)=\operatorname{loc}(\rho')\), then
\(\operatorname{st}(\rho)=\operatorname{st}_{\mathcal P}(\operatorname{loc}\rho)
=\operatorname{st}_{\mathcal P}(\operatorname{loc}\rho')=\operatorname{st}(\rho')\);
and (RC) says exactly that every residual is the locus of some record, that is,
that no residual is unstatused.
\end{proof}

\paragraph{Nonclaims.} The theorem does not
enumerate or redefine the inherited status families; that belongs to
the admissibility apparatus.  Formedness is not used: the conclusions
hold for any locus and status maps satisfying (RC) and status agreement.

\subsection*{Layer instantiations}\label{layer-inst:F9}

\begingroup\small
\begin{description}
\item[\emph{{GHS} chemical hazard classification (regulatory
registry).}] \textit{Analogy.}
GHS classification records the applicable hazard category or subcategory
using class-specific criteria, tiered assessment and weight of evidence.
Conflicting data are resolved by the applicable assessment rules. For a declared
registry, take \(R\) to be the hazard endpoints under review, \(\mathrm{Rec}\)
its endpoint-status records, \(\operatorname{loc}(\rho)\) the endpoint assessed and
\(\operatorname{st}(\rho)\) its recorded assessment outcome, including a hazard
category or an explicit nonclassification or insufficient-data status.
\Fref{F9} applies when the registry records every endpoint, satisfying (RC),
and all records at a common endpoint agree. Coverage and agreement are
conditions on this registry. EU CLP and
OSHA HazCom implement GHS-based classification requirements
\citep{UN2025GHSRev11}.

\end{description}
\endgroup


\subsection*{Stability at a glance}\label{sec:stability:summary}\addcontentsline{toc}{subsection}{Stability at a glance}
\begin{table}[H]
  \centering
  \footnotesize
  \begin{tabularx}{\textwidth}{@{}
      >{\raggedright\arraybackslash}p{0.24\textwidth}
      >{\raggedright\arraybackslash}X
      >{\raggedright\arraybackslash}p{0.12\textwidth}@{}}
    \toprule
    Name & One-line claim & Substrate cite \\
    \midrule
    No-Needles Theorem (\Fref{F6}) &
    Under the block symmetry conditions, the projection identity
    \(K_{LL}K_{LL}^{\dagger}K_{LD}=K_{LD}\) and its three budgets, no transported
    dissolving direction exceeds its declared budget (no needle).  &
    \PaperNeedleKiller, \PaperNS \\
    Sufficiency Closure Theorem (\Fref{F7}) &
    For surjective \(q\), \(q\) is sufficient for declared futures iff each future
    probe descends through \(q\).  &
    \textemdash \\
    Strict Extension Lemma (\Fref{F8}) &
    For surjective \(q_{\mathrm{old}}\), an extension is strict iff it
    retains the old quotient and splits an old fiber.  &
    \PaperAR \\
    No-Unstatused-Residual Theorem (\Fref{F9}) &
    Under residual coverage and agreement of statuses at a common locus,
    the typed status is a well-defined map on residuals and no residual is
    unstatused; under coverage, status agreement is equivalent to a
    well-defined status map, and coverage is equivalent to no unstatused
    residual.  &
    \textemdash \\
    \bottomrule
  \end{tabularx}
  \caption{Stability at a glance: the four laws of \Cref{sec:stability}.}
  \label{tab:axis-stability-summary}
\end{table}

\FloatBarrier
